% !TeX program = xelatex
\documentclass{article}
\usepackage{kotex}
\usepackage[a4paper,left=3cm, right=3cm, top=2cm, bottom=2cm]{geometry}
\usepackage{amsmath}
\usepackage{graphicx}
\usepackage{caption}
\usepackage{setspace}
\usepackage{xcolor}
\usepackage{titlesec}
\usepackage{amssymb}
\usepackage{tcolorbox}
\usepackage{wrapfig}
\usepackage{amsthm}
\usepackage{multirow}
\usepackage{array}
\usepackage{tikz}
\usetikzlibrary{positioning,arrows.meta}
\usepackage{longtable}
\usepackage{pgfplots}
\pgfplotsset{compat=1.18}
\usepackage{booktabs}

\title{2.3 Properties of Determinants; Cramer's Rule}
\date{}
\author{}
\setstretch{1.3}

\titleformat{name=\section, numberless}
  {\normalfont\large\bfseries\color{blue}}
  {}
  {0pt}
  {}
\geometry{a4paper, margin=1in}

\newcommand{\redheading}[1]{{\vspace{2mm}\noindent\Large\bfseries\color{red!55!black} #1\par}\vspace{2mm}}

\newtheoremstyle{mystyle}
  {}{}{\itshape}{}{\bfseries}{.}{.5em}{}
\theoremstyle{mystyle}

\begin{document}
\maketitle

{\itshape\color{blue!45!black}
In this section we develop fundamental properties of determinants and use them to derive
a formula for the inverse of an invertible matrix and formulas for solving certain linear systems.\par}

\vspace{2mm}

\redheading{Basic Properties of Determinants}

For $n\times n$ matrices $A$ and $B$ and any scalar $k$, we consider possible relationships
among $\det(kA)$, $\det(A+B)$, and $\det(AB)$.

Since a common factor from any row can be moved through the determinant sign, and $kA$
has $k$ as a common factor in each of its $n$ rows:
\begin{equation}
\boxed{\det(kA)=k^n\det(A)} \tag{1}
\end{equation}
For example:
\[
\begin{vmatrix}ka_{11}&ka_{12}&ka_{13}\\ka_{21}&ka_{22}&ka_{23}\\ka_{31}&ka_{32}&ka_{33}\end{vmatrix}
=k^3\begin{vmatrix}a_{11}&a_{12}&a_{13}\\a_{21}&a_{22}&a_{23}\\a_{31}&a_{32}&a_{33}\end{vmatrix}
\]
No simple relationship exists among $\det(A)$, $\det(B)$, and $\det(A+B)$; in particular,
$\det(A+B)$ usually does \emph{not} equal $\det(A)+\det(B)$.

\subsection*{EXAMPLE 1 | $\det(A+B)\neq\det(A)+\det(B)$}

Let $A=\begin{bmatrix}1&2\\2&5\end{bmatrix}$, $B=\begin{bmatrix}3&1\\1&3\end{bmatrix}$, so
$A+B=\begin{bmatrix}4&3\\3&8\end{bmatrix}$.
Then $\det(A)=1$, $\det(B)=8$, $\det(A+B)=23$; thus $\det(A+B)\neq\det(A)+\det(B)$.

There \emph{is} a useful relationship when matrices differ in exactly one row or column:

\begin{tcolorbox}[colback=red!4, colframe=red!60!black,
  title={\textbf{Theorem 2.3.1}}, boxrule=0.5mm, arc=3mm]
Let $A$, $B$, and $C$ be $n\times n$ matrices that differ only in a single row (say the
$r$th), with the $r$th row of $C$ equal to the sum of the corresponding entries of
$A$ and $B$. Then
\[
\det(C)=\det(A)+\det(B)
\]
The same result holds for columns.
\end{tcolorbox}

\subsection*{EXAMPLE 2 | Sums of Determinants}

Verify the following equality by evaluating each determinant:
\[
\det\begin{bmatrix}1&7&5\\2&0&3\\1+0&4+1&7+(-1)\end{bmatrix}
=\det\begin{bmatrix}1&7&5\\2&0&3\\1&4&7\end{bmatrix}
+\det\begin{bmatrix}1&7&5\\2&0&3\\0&1&-1\end{bmatrix}
\]

\redheading{Determinant of a Matrix Product}

Despite the complexity of matrix multiplication, a surprisingly simple relationship holds:
\begin{equation}
\det(AB)=\det(A)\det(B) \tag{2}
\end{equation}
We prove this by first establishing the following lemma.

\begin{tcolorbox}[colback=red!4, colframe=red!60!black,
  title={\textbf{Lemma 2.3.2}}, boxrule=0.5mm, arc=3mm]
If $B$ is an $n\times n$ matrix and $E$ is an $n\times n$ elementary matrix, then
\[
\det(EB)=\det(E)\det(B)
\]
\end{tcolorbox}

\textit{Proof.} Three cases corresponding to the three ERO types:

\textbf{Case 1.} $E$ results from multiplying row $i$ of $I_n$ by $k$. Then $EB$ has
row $i$ of $B$ scaled by $k$, so $\det(EB)=k\det(B)$ (Theorem 2.2.3(a)).
Also $\det(E)=k$ (Theorem 2.2.4(a)), so $\det(EB)=\det(E)\det(B)$.

\textbf{Cases 2 and 3.} The proofs for the row-interchange and row-addition cases
follow the same pattern and are left as exercises. $\blacksquare$

\textbf{Remark.} By repeated application of Lemma 2.3.2, if $E_1,E_2,\ldots,E_r$
are $n\times n$ elementary matrices and $B$ is $n\times n$:
\begin{equation}
\det(E_1 E_2\cdots E_r B)=\det(E_1)\det(E_2)\cdots\det(E_r)\det(B) \tag{3}
\end{equation}

\redheading{Determinant Test for Invertibility}

\begin{tcolorbox}[colback=red!4, colframe=red!60!black,
  title={\textbf{Theorem 2.3.3}}, boxrule=0.5mm, arc=3mm]
A square matrix $A$ is invertible if and only if $\det(A)\neq 0$.
\end{tcolorbox}

\textit{Proof.} Let $R$ be the reduced row echelon form of $A$, and let
$E_1,\ldots,E_r$ be the elementary matrices such that $R=E_r\cdots E_2 E_1 A$.
From (3):
\begin{equation}
\det(R)=\det(E_r)\cdots\det(E_2)\det(E_1)\det(A) \tag{4}
\end{equation}
Since $\det(E_i)\neq 0$ for each $i$, $\det(A)$ and $\det(R)$ are either both zero or
both nonzero.

\begin{itemize}
  \item If $A$ is invertible $\Rightarrow R=I$ (Thm.\ 1.6.4) $\Rightarrow \det(R)=1\neq 0$
    $\Rightarrow \det(A)\neq 0$.
  \item If $\det(A)\neq 0$ $\Rightarrow \det(R)\neq 0$ $\Rightarrow R$ has no zero row
    (Thm.\ 1.4.3) $\Rightarrow R=I$ $\Rightarrow A$ invertible (Thm.\ 1.6.4). $\blacksquare$
\end{itemize}

\subsection*{EXAMPLE 3 | Determinant Test for Invertibility}

Since the first and third rows of $A=\begin{bmatrix}1&2&3\\1&0&1\\2&4&6\end{bmatrix}$
are proportional, $\det(A)=0$. Thus $A$ is not invertible.

\begin{tcolorbox}[colback=red!4, colframe=red!60!black,
  title={\textbf{Theorem 2.3.4}}, boxrule=0.5mm, arc=3mm]
If $A$ and $B$ are square matrices of the same size, then
\[
\det(AB)=\det(A)\det(B)
\]
\end{tcolorbox}

\textit{Proof.}
\textbf{Case 1.} $A$ not invertible $\Rightarrow AB$ not invertible (Thm.\ 1.6.5)
$\Rightarrow \det(AB)=0$ and $\det(A)=0$, so $\det(AB)=\det(A)\det(B)=0$.

\textbf{Case 2.} $A$ invertible $\Rightarrow A=E_1 E_2\cdots E_r$ (Thm.\ 1.6.4), so
$AB=E_1\cdots E_r B$. Applying (3):
\[
\det(AB)=\det(E_1)\cdots\det(E_r)\det(B)=\det(E_1 E_2\cdots E_r)\det(B)=\det(A)\det(B). \;\blacksquare
\]

\subsection*{EXAMPLE 4 | Verifying $\det(AB)=\det(A)\det(B)$}

Let $A=\begin{bmatrix}3&1\\2&1\end{bmatrix}$, $B=\begin{bmatrix}-1&3\\5&8\end{bmatrix}$,
$AB=\begin{bmatrix}2&17\\3&14\end{bmatrix}$. Then
\[
\det(A)=1,\quad\det(B)=-23,\quad\det(AB)=-23=\det(A)\det(B).\quad\checkmark
\]

\begin{tcolorbox}[colback=red!4, colframe=red!60!black,
  title={\textbf{Theorem 2.3.5}}, boxrule=0.5mm, arc=3mm]
If $A$ is invertible, then $\displaystyle\det(A^{-1})=\frac{1}{\det(A)}$.
\end{tcolorbox}

\textit{Proof.} Since $A^{-1}A=I$, we have $\det(A^{-1})\det(A)=\det(I)=1$, so
$\det(A^{-1})=1/\det(A)$. $\blacksquare$

\redheading{Adjoint of a Matrix}

In the cofactor expansion of $\det(A)$ along row $i$, multiplying entries of row $i$
by the cofactors of a \emph{different} row $j\neq i$ always yields zero (and likewise
for columns).

\subsection*{EXAMPLE 5 | Entries and Cofactors from Different Rows}

Let $A=\begin{bmatrix}3&2&-1\\1&6&3\\2&-4&0\end{bmatrix}$.
The cofactors of $A$ are:
\[
C_{11}=12,\quad C_{12}=6,\quad C_{13}=-16
\]
\[
C_{21}=4,\quad C_{22}=2,\quad C_{23}=16
\]
\[
C_{31}=12,\quad C_{32}=-10,\quad C_{33}=16
\]

Cofactor expansion along row 1: $\det(A)=3(12)+2(6)+(-1)(-16)=36+12+16=64$.\\
Cofactor expansion along column 1: $\det(A)=3(12)+1(4)+2(12)=64$. \checkmark\\
Row 1 entries $\times$ row 2 cofactors: $3C_{21}+2C_{22}+(-1)C_{23}=12+4-16=0$.\\
Column 1 entries $\times$ column 2 cofactors: $3C_{12}+1C_{22}+2C_{32}=18+2-20=0$.

\begin{tcolorbox}[colback=teal!4, colframe=teal!65!black,
  title={\textbf{Definition 1}}, boxrule=0.5mm, arc=3mm]
If $A$ is any $n\times n$ matrix and $C_{ij}$ is the cofactor of $a_{ij}$, then the
\textbf{matrix of cofactors} from $A$ is
\[
\begin{bmatrix}C_{11}&C_{12}&\cdots&C_{1n}\\C_{21}&C_{22}&\cdots&C_{2n}\\
\vdots&\vdots&&\vdots\\C_{n1}&C_{n2}&\cdots&C_{nn}\end{bmatrix}
\]
The transpose of this matrix is called the \textbf{adjoint} of $A$, denoted $\operatorname{adj}(A)$.
\end{tcolorbox}

\subsection*{EXAMPLE 6 | Adjoint of a $3\times 3$ Matrix}

For $A=\begin{bmatrix}3&2&-1\\1&6&3\\2&-4&0\end{bmatrix}$ (from Example 5), the matrix of
cofactors is $\begin{bmatrix}12&6&-16\\4&2&16\\12&-10&16\end{bmatrix}$, so
\[
\operatorname{adj}(A)=\begin{bmatrix}12&4&12\\6&2&-10\\-16&16&16\end{bmatrix}
\]

\begin{tcolorbox}[colback=red!4, colframe=red!60!black,
  title={\textbf{Theorem 2.3.6 --- Inverse of a Matrix Using Its Adjoint}}, boxrule=0.5mm, arc=3mm]
If $A$ is an invertible $n\times n$ matrix, then
\begin{equation}
A^{-1}=\frac{1}{\det(A)}\operatorname{adj}(A) \tag{6}
\end{equation}
\end{tcolorbox}

\textit{Proof.} The $(i,j)$ entry of $A\operatorname{adj}(A)$ is:
\begin{equation}
a_{i1}C_{j1}+a_{i2}C_{j2}+\cdots+a_{in}C_{jn} \tag{7}
\end{equation}
If $i=j$: this is the cofactor expansion of $\det(A)$ along row $i$.
If $i\neq j$: entries from row $i$ multiply cofactors from row $j$, giving $0$ (Example 5).
Therefore:
\begin{equation}
A\operatorname{adj}(A)=\begin{bmatrix}\det(A)&0&\cdots&0\\0&\det(A)&\cdots&0\\
\vdots&\vdots&\ddots&\vdots\\0&0&\cdots&\det(A)\end{bmatrix}=\det(A)\cdot I \tag{8}
\end{equation}
Since $\det(A)\neq 0$, we have $A\!\left[\tfrac{1}{\det(A)}\operatorname{adj}(A)\right]=I$,
so $A^{-1}=\tfrac{1}{\det(A)}\operatorname{adj}(A)$. $\blacksquare$

\subsection*{EXAMPLE 7 | Using the Adjoint to Find an Inverse Matrix}

Use Formula (6) to find the inverse of $A$ from Example 6.

\textbf{SOLUTION:} From Example 5, $\det(A)=64$. Thus:
\[
A^{-1}=\frac{1}{\det(A)}\operatorname{adj}(A)=\frac{1}{64}\begin{bmatrix}12&4&12\\6&2&-10\\-16&16&16\end{bmatrix}
=\begin{bmatrix}\dfrac{12}{64}&\dfrac{4}{64}&\dfrac{12}{64}\\[8pt]
\dfrac{6}{64}&\dfrac{2}{64}&-\dfrac{10}{64}\\[8pt]
-\dfrac{16}{64}&\dfrac{16}{64}&\dfrac{16}{64}\end{bmatrix}
\]

\redheading{Cramer's Rule}

\begin{tcolorbox}[colback=red!4, colframe=red!60!black,
  title={\textbf{Theorem 2.3.7 --- Cramer's Rule}}, boxrule=0.5mm, arc=3mm]
If $A\mathbf{x}=\mathbf{b}$ is a system of $n$ equations in $n$ unknowns with $\det(A)\neq 0$,
the unique solution is
\[
x_1=\frac{\det(A_1)}{\det(A)},\quad
x_2=\frac{\det(A_2)}{\det(A)},\quad\ldots,\quad
x_n=\frac{\det(A_n)}{\det(A)}
\]
where $A_j$ is the matrix obtained by replacing the $j$th column of $A$ by
$\mathbf{b}=\begin{bmatrix}b_1\\b_2\\\vdots\\b_n\end{bmatrix}$.
\end{tcolorbox}

\textit{Proof.} Since $\det(A)\neq 0$, $A$ is invertible and
$\mathbf{x}=A^{-1}\mathbf{b}=\frac{1}{\det(A)}\operatorname{adj}(A)\mathbf{b}$.
The $j$th entry is:
\begin{equation}
x_j=\frac{b_1 C_{1j}+b_2 C_{2j}+\cdots+b_n C_{nj}}{\det(A)} \tag{9}
\end{equation}
Since $A_j$ differs from $A$ only in column $j$, the cofactors of that column in $A_j$
equal $C_{1j},\ldots,C_{nj}$ of $A$. Hence $\det(A_j)=b_1 C_{1j}+\cdots+b_n C_{nj}$,
giving $x_j=\det(A_j)/\det(A)$. $\blacksquare$

\subsection*{EXAMPLE 8 | Using Cramer's Rule to Solve a Linear System}

Use Cramer's rule to solve:
\[
x_1+2x_3=6,\quad -3x_1+4x_2+6x_3=30,\quad -x_1-2x_2+3x_3=8
\]

\textbf{SOLUTION:}
\[
A=\begin{bmatrix}1&0&2\\-3&4&6\\-1&-2&3\end{bmatrix},\quad
A_1=\begin{bmatrix}6&0&2\\30&4&6\\8&-2&3\end{bmatrix},\quad
A_2=\begin{bmatrix}1&6&2\\-3&30&6\\-1&8&3\end{bmatrix},\quad
A_3=\begin{bmatrix}1&0&6\\-3&4&30\\-1&-2&8\end{bmatrix}
\]
Therefore:
\[
x_1=\frac{\det(A_1)}{\det(A)}=\frac{-40}{44}=-\frac{10}{11},\qquad
x_2=\frac{\det(A_2)}{\det(A)}=\frac{72}{44}=\frac{18}{11},\qquad
x_3=\frac{\det(A_3)}{\det(A)}=\frac{152}{44}=\frac{38}{11}
\]

\redheading{Equivalence Theorem}

By merging Theorem 2.3.3 with the earlier Theorem 1.6.4, we obtain:

\begin{tcolorbox}[colback=red!4, colframe=red!60!black,
  title={\textbf{Theorem 2.3.8 --- Equivalent Statements}}, boxrule=0.5mm, arc=3mm]
If $A$ is an $n\times n$ matrix, the following statements are equivalent:
\begin{enumerate}
  \item[(a)] $A$ is invertible.
  \item[(b)] $A\mathbf{x}=\mathbf{0}$ has only the trivial solution.
  \item[(c)] The reduced row echelon form of $A$ is $I_n$.
  \item[(d)] $A$ can be expressed as a product of elementary matrices.
  \item[(e)] $A\mathbf{x}=\mathbf{b}$ is consistent for every $n\times 1$ matrix $\mathbf{b}$.
  \item[(f)] $A\mathbf{x}=\mathbf{b}$ has exactly one solution for every $n\times 1$ matrix $\mathbf{b}$.
  \item[(g)] $\det(A)\neq 0$.
\end{enumerate}
\end{tcolorbox}

\vspace{5mm}

\noindent\rule{\linewidth}{0.5pt}

\section*{Exercise Set 2.3}

\noindent\textit{In Exercises 1--3, verify that $\det(kA)=k^n\det(A)$.}

\begin{enumerate}

\item[1.] $A=\begin{bmatrix}-1&2\\3&4\end{bmatrix}$;\; $k=2$

\item[3.] $A=\begin{bmatrix}2&-1&3\\3&2&1\\1&4&5\end{bmatrix}$;\; $k=-2$

\noindent\textit{In Exercise 5, verify that $\det(AB)=\det(BA)$ and determine whether
$\det(A+B)=\det(A)+\det(B)$.}

\item[5.] $A=\begin{bmatrix}2&1&0\\3&4&0\\0&0&2\end{bmatrix}$,\quad
$B=\begin{bmatrix}1&-1&3\\7&1&2\\5&0&1\end{bmatrix}$

\noindent\textit{In Exercises 7--14, use determinants to decide whether the matrix is
invertible.}

\item[7.] $A=\begin{bmatrix}2&5&5\\-1&-1&0\\2&4&3\end{bmatrix}$

\item[9.] $A=\begin{bmatrix}2&-3&5\\0&1&-3\\0&0&2\end{bmatrix}$

\item[11.] $A=\begin{bmatrix}4&2&8\\-2&1&-4\\3&1&6\end{bmatrix}$

\item[14.] $A=\begin{bmatrix}\sqrt{2}&-\sqrt{7}&0\\3\sqrt{2}&-3\sqrt{7}&0\\5&-9&0\end{bmatrix}$

\noindent\textit{In Exercises 15--17, find the values of $k$ for which the matrix $A$
is invertible.}

\item[15.] $A=\begin{bmatrix}k-3&-2\\-2&k-2\end{bmatrix}$

\item[17.] $A=\begin{bmatrix}1&2&4\\3&1&6\\k&3&2\end{bmatrix}$

\noindent\textit{In Exercises 19 and 22, decide whether the matrix is invertible; if so,
use the adjoint method to find its inverse.}

\item[19.] $A=\begin{bmatrix}2&5&5\\-1&-1&0\\2&4&3\end{bmatrix}$

\item[22.] $A=\begin{bmatrix}2&0&0\\8&1&0\\-5&3&6\end{bmatrix}$

\noindent\textit{In Exercises 24--27, solve by Cramer's rule, where it applies.}

\item[24.] $7x_1-2x_2=3$,\quad $3x_1+x_2=5$

\item[25.] $4x+5y=2$,\quad $11x+y+2z=3$,\quad $x+5y+2z=1$

\item[27.] $x_1-3x_2+x_3=4$,\quad $2x_1-x_2=-2$,\quad $4x_1-3x_3=0$

\item[30.] Show that
$A=\begin{bmatrix}\cos\theta&\sin\theta&0\\-\sin\theta&\cos\theta&0\\0&0&1\end{bmatrix}$
is invertible for all values of $\theta$; then find $A^{-1}$ using Theorem~2.3.6.

\item[33.] Let $A=\begin{bmatrix}a&b&c\\d&e&f\\g&h&i\end{bmatrix}$ with $\det(A)=-7$.
Find:
\begin{enumerate}
  \item[(a)] $\det(3A)$
  \item[(b)] $\det(A^{-1})$
  \item[(c)] $\det(2A^{-1})$
  \item[(d)] $\det\bigl((2A)^{-1}\bigr)$
  \item[(e)] $\det\begin{bmatrix}a&g&d\\b&h&e\\c&i&f\end{bmatrix}$
\end{enumerate}

\item[34.] In each part, find the determinant given that $A$ is a $4\times 4$ matrix
with $\det(A)=-2$:
\begin{enumerate}
  \item[(a)] $\det(-A)$
  \item[(b)] $\det(A^{-1})$
  \item[(c)] $\det(2A^T)$
  \item[(d)] $\det(A^3)$
\end{enumerate}

\end{enumerate}

\end{document}
